\section{Related Work}\label{sec:related}

\subsection{Feature Selection Paradigms}

Feature selection is a fundamental pre-processing step in high-dimensional data analysis, aiming to identify a minimal subset of relevant features that preserves or improves classification performance while reducing computational cost and enhancing interpretability \cite{lili2017,guyon2003}. Methods fall broadly into three paradigms: filter, wrapper, and embedded.

Filter methods evaluate features independently of any learning algorithm, relying on intrinsic data properties such as variance, correlation, or eigenvector centrality \cite{kramer1991,roffo2016}. Representative techniques include the minimum Redundancy Maximum Relevance (mRMR) criterion \cite{peng2005} and the Fast Correlation-Based Filter (FCBF) \cite{yu2003}. These methods are computationally efficient and scale readily to high-dimensional settings, yet they disregard feature interactions and may retain redundant subsets under strong inter-feature correlation.

Wrapper methods treat feature selection as a search problem, evaluating candidate subsets by the predictive performance of a specific classifier \cite{kohavi1997}. Hybrid genetic algorithms \cite{oh2004} and other meta-heuristic wrappers can capture complex feature interdependencies, but their computational cost grows rapidly with dimensionality, as each candidate subset requires a full classifier training-and-evaluation cycle.

Embedded methods integrate feature selection into model training. Lasso and its bootstrapped variant Bolasso \cite{bach2008} enforce sparsity through $L_1$ regularization; gradient-boosting-based approaches such as FS-GBDT \cite{zhang2021fsgbdt} leverage tree-split importance; and hybrid algorithms like ILRC \cite{yu2021ilrc} combine regularization with clustering. Embedded methods offer a favorable speed--accuracy trade-off, but their feature rankings are tied to specific model assumptions.

At high dimensionality, each paradigm faces distinct structural bottlenecks: filters are fast but oblivious to feature interactions, wrappers are accurate but face a search space that grows exponentially with dimension, and embedded methods depend on model-specific regularization biases. These complementary weaknesses motivate evolutionary computation, which decouples the search mechanism from the evaluation metric and navigates the exploration--exploitation trade-off through population-based global search.

\subsection{Single-Objective Evolutionary Feature Selection}

Evolutionary algorithms have been extensively applied to feature selection for their population-based global search and gradient-free nature. Early work adapted classical meta-heuristics---binary PSO (BPSO) \cite{kennedy1997}, ant colony optimization (ACO) \cite{aghdam2009}, artificial bee colony (ABC) \cite{karaboga2007}, and genetic algorithms with support vector machines (GA-SVM) \cite{huang2006}---to the binary subset selection problem, establishing the template of a population of candidate subsets evolving under selection, crossover, and mutation guided by a scalar fitness function.

Because standard evolutionary operators degrade at high dimensionality, specialized single-objective methods have been developed: self-adaptive PSO adjusting swarm size and mutation rate \cite{xue2019}, progressively refined PSO alternating between coarse- and fine-grained search \cite{prasad2015}, and problem-specific heuristics---FWPSO~\cite{wangjia2022}, which biases velocity updates with per-feature importance; SaWDE~\cite{wang2022} with adaptive weighting; IGPSO~\cite{xuezhang2024}, which estimates relevance via a multi-layer perceptron; and PSO-CSM~\cite{wei2025} integrating multiple evaluation criteria. Learning-based control and space decomposition are also explored: reinforcement-learning-aided PSO~\cite{weil2025} and divide-and-conquer search~\cite{chen2025}. These methods refine search operators but do not systematically exploit accumulated feature-importance knowledge harvested from the search itself.

\subsection{Multi-Objective Evolutionary Feature Selection}

Multi-objective formulations treat classification accuracy and feature count as competing objectives, producing a Pareto front of non-dominated solutions. Xue et al. \cite{xue2012mopso} pioneered multi-objective PSO for feature selection, showing that jointly optimizing accuracy and subset size yields more diverse and informative solutions than scalarized approaches. Hu et al. \cite{hu2020} extended this framework with fuzzy cost functions for imbalanced classification.

The multi-objective evolutionary algorithm (MOEA) family provides several architectures for this problem. NSGA-II-FS \cite{deb2002} uses non-dominated sorting and crowding distance to maintain Pareto-front diversity. MOEA/D-FS \cite{zhang2007} decomposes the problem into scalar subproblems solved cooperatively. MOPSO-FS \cite{coello2004} extends single-objective PSO with an external archive and Pareto-dominance-based leader selection. For high-dimensional data, Cheng et al. \cite{cheng2023} proposed a variable-granularity search method (VGS-MOEA) that transitions from coarse to fine-grained search, and Hancer et al. \cite{hancer2023} combined evolutionary search with filter-based ranking under both single- and multi-objective settings.

Despite their conceptual advantages, MOFS methods incur substantial overhead: maintaining an external archive, non-dominated sorting, and Pareto-front management all cost significant time (in our experiments, over 110 seconds on average). Moreover, under fixed-length binary encoding at high dimensionality, the Pareto front expands combinatorially and convergence slows, limiting scalability at gene-expression scale.

\subsection{Evolutionary Multi-Task Learning}

Evolutionary multi-task learning (EMTL) solves multiple related tasks jointly, with knowledge transfer accelerating convergence. Zhang and Yang \cite{zhang2021mtl} surveyed its foundations, and several instantiations target feature selection: a minion-based multitasking framework \cite{zhang2021minion}, PSO-based multitasking for high-dimensional classification \cite{chen2020emfs,chen2021emfs}, multitask bee colony optimization for hyperspectral band selection \cite{he2022}, and multitasking for large-scale multi-objective optimization \cite{liu2022}.

Of particular relevance, Wang et al. \cite{mel2024} proposed MEL (Multi-task Evolutionary Learning), a dual-swarm multi-task PSO in which one swarm optimizes binary feature subsets while a second learns continuous feature weights, with the swarms cooperating through cross-task knowledge transfer. MEL was selected as the backbone for three reasons: its fixed-length encoding is efficient at high dimensionality and naturally compatible with incremental augmentation; its dual-space cooperation provides a built-in mechanism for capturing feature relevance signals; and it has demonstrated strong performance on gene-expression benchmarks.

Nevertheless, MEL exhibits three structural limitations. Its population is initialized entirely at random, wasting search budget in the sparse high-dimensional space. Its operators---velocity update and mutation---are uniform and stochastic, never incorporating cross-iteration knowledge about which features have proven beneficial. And its dual-objective weights remain fixed throughout search, imposing a rigid schedule that cannot adapt to the natural phasing of feature selection, where early iterations benefit from aggressive compression and later iterations demand fine-grained accuracy refinement.

\subsection{Gene-Expression Feature Selection}

High-dimensional gene expression data present a particularly challenging testbed: the number of features routinely exceeds the number of samples by three to four orders of magnitude---a ``small $N$, large $D$'' regime where conventional statistical methods break down and overfitting is pervasive \cite{bermingham2015,phuong2005}. Representative applications include cancer staging \cite{saghapour2017}, Alzheimer's classification from neuroimaging features \cite{zhang2015eigenbrain}, biomarker gene identification \cite{wangjia2022}, and multi-cancer risk module discovery \cite{zhang2021fsgbdt,yu2021ilrc}.

Two structural properties amplify the difficulty. First, feature redundancy is exceptionally high: co-regulated genes exhibit correlated expression, so distinct subsets frequently yield near-identical accuracy. Second, a non-trivial fraction of gene datasets reach accuracy ceilings---exceeding 0.998 under multiple methods---rendering accuracy alone insufficient to distinguish algorithmic quality; on these datasets the selected feature count becomes the primary discriminatory metric. Evaluation must therefore jointly consider classification accuracy and subset compactness.

Recent methodological development in this area is extensive; surveys of evolutionary computation for biomarker identification in complex diseases \cite{liang2025} and instance-selection-assisted evolutionary feature selection \cite{xia2026} both target the small-sample, ultra-high-dimensional regime central to gene-expression analysis.

\subsection{Positioning of KDS-EL}

KDS-EL builds upon the multi-task PSO backbone of MEL, adding three mechanisms that each target one structural limitation: MI-guided initialization with adaptive re-initialization addresses the cold-start problem; the KER strategy injects cross-iteration feature importance into elite refinement at no additional cost; and Two-Phase Weight Scheduling replaces static weights with a compression-then-refinement schedule. These mechanisms preserve the fixed-length encoding that makes the backbone suitable for high-dimensional applications.
